\section*{Note on version 1.1}
Version 1.1 corrects two statements of version 1.0 and adds the states behind the corrections. (i) The smallest gap found for the pendulum is $0.916\,\sigma$, not $0.946\,\sigma$; the numerical estimate of the constant in the resolution-gap conjecture is lowered accordingly, and on the new states the state-wise bound of Theorem~\ref{thm:statewise}, which holds on all of them, is at least $0.869\,\sigma$ (Secs.~\ref{sec:bound} and~\ref{sec:evidence}). (ii) The ``bent needle'' of the twist flow is a chain of straight needles whose gap depends on the number of pieces; its agreement with the value $1/2$ is not established for a smooth needle (Numerical result~\ref{obs:bent}). The code and the verification suite were extended accordingly. No theorem is affected.

\section*{Acknowledgments}
The author used two AI assistants in this work: Google Gemini (\mbox{Gemini 3.8 Flash}), for independent cross-checking of the derivations and results from several angles; and Anthropic Claude (\mbox{claude-opus-5-5}), in two separate sessions. In the first session it took part in theoretical discussion, proposed and checked proofs, derived the one-variable form~\eqref{eq:Rstar}, wrote independent numerical checks, and drafted and edited the text of the manuscript. In the second session it designed and ran the numerical study, wrote the evaluation and verification code, derived Theorem~\ref{thm:statewise}, Proposition~\ref{prop:smooth}, the closed form~\eqref{eq:quadratic} and the series~\eqref{eq:series}, and generated the figures. Results were passed between the sessions as written summaries, and every summary was checked in the first session by re-running the full verification suite and, for the central numbers, by the independent methods of Appendix~\ref{app:methods}. The instructions given to the assistants consisted of the author's research questions and requests to derive, check, critique, compute or rephrase specific results. Several turns of the study were proposed by the author: to observe the nonlinear problem numerically before attempting a proof, to examine the spatial spread of the local anisotropy, which identified the covariance term as the mechanism of excess over single components, to ask whether the flow singles out a reference frame, which led to the formulation in the stretching and contracting directions, to follow the rotation of that frame in flows where it turns, and to report the refutation of the extremality of Gaussian states as a result in its own right. No result was adopted without verification: every analytical statement, number and figure was checked against independent numerical evaluation in the openly available code, and all bibliographic information was checked against the original sources. The author directed the research, made all decisions on its content, and takes full responsibility for this article.

\section*{Data and Code Availability}
Scripts that reproduce every analytical statement, number and figure of this article, together with the verification suite, are openly available at \url{https://github.com/neginyan/resolution-gap-arrow-of-time}. The version used for this article is archived at Zenodo, \href{https://doi.org/10.5281/zenodo.23318980}{doi:10.5281/zenodo.23318980}.
